\documentclass[12pt]{article}
\usepackage{amsmath}
\usepackage{enumitem}
\begin{document}

\section{Ghost Trace Test}

This document tests whether incremental dirty rect rendering leaves ghost traces
after editing. Open the built-in preview, then perform the edits below.

\subsection{Test 1: Delete a middle list item}

\begin{enumerate}
      \item First item with enough text to wrap across multiple lines so the paragraph is
            wide and fills a substantial portion of the page width for testing dirty rect
            coverage when content is deleted and surrounding text reflows upward.
      \item \textbf{DELETE THIS ENTIRE LINE} -- second item that also contains substantial
            text wrapping across multiple lines at a different horizontal position than the
            items above and below it in the list for triggering the dirty rect edge-case bug.
      \item Third item with more content filling several lines to ensure we have enough
            vertical space for the dirty rect algorithm to encounter rows with varying
            horizontal pixel positions after the middle item is removed from the document.
\end{enumerate}

\subsection{Test 2: Delete a blockquote}

Normal paragraph at full page width that is positioned above the blockquote.
When the quote below is deleted, this text stays in place but the text below
it will shift upward, creating a dirty region that spans different horizontal
positions.

\begin{quote}
      This blockquote is indented from both sides. When this entire block is deleted,
      the paragraph below will move up into this position. The dirty rect must cover
      the full width of the quote text AND the moved paragraph text, which start at
      different horizontal offsets -- triggering the edge-case bug if the rect width
      is computed incorrectly in the single-pass loop.
\end{quote}

Another normal paragraph at full page width. When the blockquote above is
deleted, this text shifts upward. The area where the blockquote was (indented)
is now filled with this full-width text. The dirty rect algorithm must handle
the varying min\_x values across rows correctly.

\subsection{Test 3: Rapid undo/redo}

Type some random characters here, then undo them. Watch for any visual artifacts
or ghost traces that remain after the undo operation completes.

Another normal paragraph at full page width. When the blockquote above is
deleted, this text shifts upward. The area where the blockquote was (indented)
is now filled with this full-width text. The dirty rect algorithm must handle
the varying min\_x values across rows correctly.

Another normal paragraph at full page width. When the blockquote above is
deleted, this text shifts upward. The area where the blockquote was (indented)
is now filled with this full-width text. The dirty rect algorithm must handle
the varying min\_x values across rows correctly.

Another normal paragraph at full page width. When the blockquote above is
deleted, this text shifts upward. The area where the blockquote was (indented)
is now filled with this full-width text. The dirty rect algorithm must handle
the varying min\_x values across rows correctly.

Another normal paragraph at full page width. When the blockquote above is
deleted, this text shifts upward. The area where the blockquote was (indented)
is now filled with this full-width text. The dirty rect algorithm must handle
the varying min\_x values across rows correctly.

Another normal paragraph at full page width. When the blockquote above is
deleted, this text shifts upward. The area where the blockquote was (indented)
is now filled with this full-width text. The dirty rect algorithm must handle
the varying min\_x values across rows correctly.

\begin{subequations}
      \begin{align}
            a & = b + c      \\
            d & = e - f      \\
            g & = h \times i \\
            j & = k \div l   \\
            p & = \sqrt{q}   \\
            r & = \log(s)    \\
            t & = \sin(u)    \\
            v & = \cos(w)    \\
      \end{align}
\end{subequations}

\end{document}
